{{دیوان
| عنوان = عبد الرّحمٰن اوّل کا بویا ہُوا کھجور کا پہلا درخت
| ماخذ = ویکی ماخذ
| ماخذ_لنک = https://ur.wikisource.org/wiki/%D8%B9%D8%A8%D8%AF_%D8%A7%D9%84%D8%B1%D9%91%D8%AD%D9%85%D9%B0%D9%86_%D8%A7%D9%88%D9%91%D9%84_%DA%A9%D8%A7_%D8%A8%D9%88%DB%8C%D8%A7_%DB%81%D9%8F%D9%88%D8%A7_%DA%A9%DA%BE%D8%AC%D9%88%D8%B1_%DA%A9%D8%A7_%D9%BE%DB%81%D9%84%D8%A7_%D8%AF%D8%B1%D8%AE%D8%AA
| مقابلہ = اقبال سائبر لائبریری (اقبال اکادمی پاکستان)
| مقابلہ_لنک = https://www.iqbalcyberlibrary.net
}}
=== سرزمینِ اندلس میں ===

یہ اشعار جو عبد الرحمٰن اوّل کی تصنیف سے ہیں، ’تاریخ المقری‘ میں درج ہیں مندرجہ ذیل اردو نظم ان کا آزاد ترجمہ ہے (درخت مذکور مدینتہ الزّہرا میں بویا گیا تھا)

<poem>
میری آنکھوں کا نُور ہے تُو
میرے دل کا سُرور ہے تُو

اپنی وادی سے دُور ہوں مَیں
میرے لیے نخلِ طُور ہے تُو

مغرب کی ہوا نے تجھ کو پالا
صحرائے عرب کی حُور ہے تُو

پردیس میں ناصبور ہوں مَیں
پردیس میں ناصبُور ہے تو

غُربت کی ہوا میں باروَر ہو
ساقی تیرا نمِ سحَر ہو

عالَم کا عجیب ہے نظارہ
دامانِ نِگہ ہے پارہ پارہ

ہمّت کو شناوری مبارک!
پیدا نہیں بحر کا کنارہ

ہے سوزِ دُروں سے زندگانی
اُٹھتا نہیں خاک سے شرارہ

صُبحِ غُربت میں اور چمکا
ٹُوٹا ہُوا شام کا ستارہ

مومن کے جہاں کی حد نہیں ہے
مومن کا مقام ہر کہیں ہے
</poem>
